% concatenated from za.tex
%% This is file `elsarticle-template-1-num.tex',
%%
%% Copyright 2009 Elsevier Ltd
%%
%% This file is part of the 'Elsarticle Bundle'.
%% ---------------------------------------------
%%
%% It may be distributed under the conditions of the LaTeX Project Public
%% License, either version 1.2 of this license or (at your option) any
%% later version.  The latest version of this license is in
%%    http://www.latex-project.org/lppl.txt
%% and version 1.2 or later is part of all distributions of LaTeX
%% version 1999/12/01 or later.
%%
%% The list of all files belonging to the 'Elsarticle Bundle' is
%% given in the file `manifest.txt'.
%%
%% Template article for Elsevier's document class `elsarticle'
%% with numbered style bibliographic references
%%
%% $Id: elsarticle-template-1-num.tex 149 2009-10-08 05:01:15Z rishi $
%% $URL: http://lenova.river-valley.com/svn/elsbst/trunk/elsarticle-template-1-num.tex $
%%
\documentclass[preprint,12pt]{elsarticle1}

%% Use the option review to obtain double line spacing
%% \documentclass[preprint,review,12pt]{elsarticle}

%% Use the options 1p,twocolumn; 3p; 3p,twocolumn; 5p; or 5p,twocolumn
%% for a journal layout:
%% \documentclass[final,1p,times]{elsarticle}
%% \documentclass[final,1p,times,twocolumn]{elsarticle}
%% \documentclass[final,3p,times]{elsarticle}
%% \documentclass[final,3p,times,twocolumn]{elsarticle}
%% \documentclass[final,5p,times]{elsarticle}
%% \documentclass[final,5p,times,twocolumn]{elsarticle}

%% if you use PostScript figures in your article
%% use the graphics package for simple commands
%% \usepackage{graphics}
%% or use the graphicx package for more complicated commands
%% \usepackage{graphicx}
%% or use the epsfig package if you prefer to use the old commands
%% \usepackage{epsfig}
%% The amssymb package provides various useful mathematical symbols
\addtolength{\topmargin}{-9mm}
\setlength{\oddsidemargin}{5mm}  % 3.17cm - 1 inch
\setlength{\evensidemargin}{0mm}
\setlength{\textwidth}{15cm}
\setlength{\textheight}{21cm}
\usepackage{subfig}
\usepackage{graphicx}
\usepackage{caption,color}   % 24.62
\usepackage{amssymb}
\usepackage{amsthm}
\usepackage{amsmath}
\usepackage{epic}
\usepackage{setspace}
\usepackage{float}
\usepackage{multirow}
%\usepackage{cite}
\newtheorem{thm}{Theorem}[section]
\newtheorem{cor}[thm]{Corollary}
\newtheorem{con}[thm]{Conjecture}
\newtheorem{lem}[thm]{Lemma}
\newtheorem{pro}[thm]{Proposition}
\newtheorem{defi}[thm]{Definition}
\newtheorem{remark}[thm]{Remark}
\newtheorem{example}[thm]{Example}
\newtheorem{state}[thm]{Statement}
\newcommand{\heiti}{\CJKfamily{hei}}
\newcommand{\ixiaosihao}{\fontsize{5pt}{\baselineskip}\selectfont}
\newcommand{\texiaosihao}{\fontsize{7pt}{\baselineskip}\selectfont}
\newcommand{\txiaosihao}{\fontsize{8pt}{\baselineskip}\selectfont}
\newcommand{\xiaosihao}{\fontsize{9pt}{\baselineskip}\selectfont}
\newcommand{\dasihao}{\fontsize{14pt}{\baselineskip}\selectfont}
\newcommand{\minitab}[2][l]{\begin{tabular}{#1}#2\end{tabular}}
\newcommand{\red}[1]{{\color{red}#1}}
\newcommand{\blue}[1]{{\color{blue}#1}}
\newenvironment{proofa}{\begin{proof}[Proof of Theorem 1.3]}{\end{proof}}
\newenvironment{proofb}{\begin{proof}[Proof of Theorem 1.2]}{\end{proof}}
\numberwithin{equation}{section}
\journal{}


\begin{document}
\begin{spacing}{1.15}
\begin{frontmatter}
\title{Extremal Zagreb indices of bicyclic hypergraphs}

\author{Hong Zhou}
\author{Changjiang Bu}\ead{buchangjiang@hrbeu.edu.cn}
\address{School of Mathematical Sciences, Harbin Engineering University, Harbin 150001, PR China}

\begin{abstract}

The Zagreb index of a hypergraph is defined as the sum of the squares of the degrees of its vertices.
A connected $k$-uniform hypergraph with $n$ vertices and $m$ edges is called bicyclic if $n=m(k-1)-1$.
In this paper, we determine the hypergraphs with the maximum and minimum Zagreb indices among all linear bicyclic uniform hypergraphs.

\end{abstract}


\begin{keyword}
bicyclic hypergraph, Zagreb index\\
\emph{MSC:}
05C65, 05C09
\end{keyword}
\end{frontmatter}



\section{Introduction}

Let $G=(V,E)$ be a simple graph with the vertex set $V(G)$ and the edge set $E(G)$.
The first Zagreb index $M_{1}$ of a graph $G$ is defined as
$$
M_{1}(G)=\sum\limits_{u\in V(G)}(d_{G}(u))^{2},
$$
where $d_{G}(u)$ is the degree of vertex $u$ in $G$. The first Zagreb index is also called Gutman index.


The first Zagreb index, first proposed by Gutman and Trinajsti{\'c} \cite{1972Graph}, is originated
from chemical studies on total $\pi$-electron energy of conjugated molecules.
The main properties of first Zagreb indices were summarized in \cite{30years,2004Graph}.
%Gutman and Das \cite{2004Graph} showed that the trees with the smallest and largest $M_{1}$ are the path and the star, respectively.
Deng \cite{Aunifiedapproach} characterized the graphs with the maximum and minimum first Zagreb indices among all bicyclic graphs.
Some results on extremal first Zagreb indices have been obtained in the literature:
see \biboptions{sort&compress}\cite{2004Graph,treedegree,treesdomination,treesdistancedomination} for trees, \cite{Uncyclicgraphswiththefirstthree,orderUncyclicgraphs} for unicyclic graphs,
\cite{Zagrebindicesquasiunicyclic}  for $k$-generalized quasi unicyclic graphs,
\cite{Zagrebindices} for triangle-free graphs, and \biboptions{sort&compress}\cite{Zagrebindicescut,Zagrebindicescutedges,Zagrebindicesclique,Zagrebindicesconnectivity,Zagrebindicespendent} for graphs with given parameters.




%see \biboptions{sort&compress}\cite{WU20101707,PAN20111265,CHENG2012858} for trees, \cite{CHENG20121123} for unicyclic graphs and  \cite{SHUCHAO2013ON,10.1216/RMJ-2016-46-1-261}



%The study of spectral hypergraph theories via tensors has attracted extensive attention \biboptions{sort&compress}\cite{clark2021harary,cooper2,doi:10.1137/21M1404740,2022Thestabilizing}, especially the spectral moment of hypergraphs \cite{clark2021harary,ref23}.

%Hypergraphs are used in chemistry to model molecules or chemical reactions involving simultaneous bonding of multiple atoms.
Hypergraphs find application in chemistry when modeling molecules or
chemical reactions involving multiple atoms bonding simultaneously.
Hypergraphs offer a more accurate depiction of certain
chemical scenarios, such as transition states in reactions, which involve
multiple atoms simultaneously changing their bonding configurations. The model for an organometallic
compound, where the hyperedges with two vertices represent covalent
bond and with more than two vertices represent delocalised polycentric
bonds, was studied in \cite{Applicationfhypergraph}.

For a hypergraph $\mathcal{H}$, $$\sum\limits_{u\in V(\mathcal{H})}(d_{\mathcal{H}}(u))^{2}$$ is called the Zagreb index of $\mathcal{H}$, denoted by $M(\mathcal{H})$, where $d_{\mathcal{H}}(u)$ is the degree of vertex $u$ in $\mathcal{H}$ \cite{Kau2020Energies}.
The hypergraphs with the maximum and minimum Zagreb indices were determined both for uniform hypertrees and for linear unicyclic uniform hypergraphs \cite{2309.16925}.
The bounds on the Zagreb indices of hypergraphs, weak bipartite hypergraphs, hypertrees, $k$-uniform
hypergraphs, $k$-uniform weak bipartite hypergraphs, and $k$-uniform hypertrees were given \cite{zagrebhy}.
%The hypergraphs with the maximum and minimum Zagreb indices among all  hypertrees were given. And the hypergraphs with the maximum and minimum Zagreb indices in unicyclic hypergraph were given


In this paper, we characterize the hypergraphs with the maximum and minimum Zagreb indices among all linear bicyclic uniform hypergraphs.



\section{Preliminaries}

A hypergraph $\mathcal{H}$ is called \textit{$k$-uniform} if every edge of $\mathcal{H}$ contains exactly $k$ vertices.
%The eigenvalues of $\mathcal{A}_\mathcal{H}$ are also called the eigenvalues of $\mathcal{H}$ .
%The \textit{degree} of a vertex $v$ of $\mathcal{H}$ is the number of edges containing the vertex, denoted by $d_{\mathcal{H}}(v)$ or $d(v)$.
A vertex of $\mathcal{H}$ is called a \textit{cored vertex} if it has degree one.
%A vertex with degree larger than one is called an \textit{intersection vertex}.
An edge $e$ of $\mathcal{H}$ is called a \textit{pendant edge} if it contains exactly $|e|-1$ cored vertices.
A cored vertex in a pendant edge is also called a \textit{pendant vertex}.
A hypergraph $\mathcal{H}$ is called \textit{linear} if any two edges intersect into at most one vertex.
The \textit{girth} of $\mathcal{H}$ is the minimum length of the hypercycles of $\mathcal{H}$.
%, denoted by $g(\mathcal{H})$.
%$\mathcal{H}$ is called \textit{trivial} if it contains only one vertex; otherwise, it is called
%\textit{nontrivial}.
%$\mathcal{H}$ is called \textit{nontrivial} if it contains more than one vertex.
%The \textit{degree} of a vertex $v$ of $\mathcal{H}$ is the number of edges containing the vertex, denoted by $d_{\mathcal{H}}(v)$ or $d(v)$. or $d(v)$
%$\mathcal{H}$ is called \textit{trivial} if it contains only one vertex.
%The \textit{$k$-power hypergraph} $G^{(k)}$ is the hypergraph which is obtained by adding $k-2$ vertices with degree one to each edge of the graph $G$.
%Note if $m=2$, $G^{(m)}=G.$
%Let $P_{q}$ be a path of length $q$, $S_{q}$ be a star with $q$ edges and $C_{n}$ be a cycle with $n$ edges.
%Let $P_{q}$, $S_{q}$ and $C_{q}$ denote the path, the star and the cycle with $q$ edges, respectively.
%A connected $k$-uniform hypergraph $\mathcal{H}$ with $n$ vertices and $m$ edges is called bicyclic if $n=m(k-1)-1$.
%There are the following two types of all linear bicyclic $k$-uniform hypergraphs.
All linear bicyclic $k$-uniform hypergraphs with $n$ vertices and $m$ edges consist of the following two types $\mathcal{B}_{n}^{k}$ and $\mathcal{C}_{n}^{k}$.


For three integers $p, q, l$ with $q\geq p\geq3$ and $l\geq 0$, let $C_{1} (resp., C_{2})$ be a $k$-uniform
 hypercycle of length $p$ (resp., $q$) and $P = (u_{0},e_{1},u_{1},\ldots,e_{l},u_{l})$ be a $k$-uniform hyperpath
of length $l$. Let $v_{1,1}\in V(C_{1}), v_{2,1}\in V(C_{2})$ be two arbitrary vertices with degree 1, and
let $v_{1,2}\in V(C_{1}), v_{2,2}\in V(C_{2})$ be two arbitrary vertices with degree 2. For $n'= (p+q+l)(k-1)-1$, let $B_{1,n',p,l,q}^{k}$
be the $n'$-vertex $k$-uniform bicyclic hypergraph obtained by
identifying $v_{1,2}$ with $u_{0}$, and identifying $v_{2,2}$ with $u_{l}$, let $B_{2,n',p,l,q}^{k}$
be the $n'$-vertex $k$-uniform
bicyclic hypergraph obtained by identifying $v_{1,2}$ with $u_{0}$, and identifying $v_{2,1}$ with $u_{l}$, and
let $B_{3,n',p,l,q}^{k}$
be the $n'$-vertex $k$-uniform bicyclic hypergraph obtained by identifying $v_{1,1}$ with $u_{0}$, and identifying $v_{2,1}$ with $u_{l}$.

For $n\geq n'$ and $i\in \{1, 2, 3\}$, let $\mathcal{B}_{i,n,p,l,q}^{k}$ be the set of $n$-vertex $m$-edge $k$-uniform bicyclic hypergraphs
each of which contains $B_{i,n',p,l,q}^{k}$ as a sub-hypergraph, where $m=\frac{n+1}{k-1}$. %And $B_{i,n',p,l,q}^{k}$ is called the base of $\mathcal{B}_{i,n,p,l,q}^{k}$.
%Let $B_{n,p,l,q}^{k}=\bigcup_{i=1}^{3}B_{i,n,p,l,q}^{k}$.
Let $\mathcal{B}_{n}^{k}=\bigcup_{i=1}^{3}\{\mathcal{B}_{i,n,p,l,q}^{k}~|~q\geq p\geq3,l\geq 0\}$.
Moreover, let $B_{1,n,p,0,q}^{k}(m-p-q)$ denote the $k$-uniform bicyclic hypergraph
obtained from $B_{1,n',p,0,q}^{k}$ by attaching $m-p-q$ pendant edges at the unique vertex with
degree 4, where $m=\frac{n+1}{k-1}$.

%For three integers $p, q, l$ with $1\leq p\leq q\leq l$ and when $p=1$, $q>p$.\geq ,1\leq p\leq \min\{q,l\},l\geq q-2,1\leq p\leq \min\{q,l\}
Let $P_{p}=(u_{1},e_{1},u_{2},\ldots,e_{p},u_{p+1}), P_{q}=(v_{1}, f_{1}, v_{2},\ldots, f_{q}, v_{q+1})$ and $P_{l}=(w_{1},g_{1},\\w_{2},\ldots,g_{l},w_{l+1})$ be three $k$-uniform hyperpaths, and suppose $(p+q+l)(k-1)-1=n'$. For three integers $p, q, l$ with $p=1,1<q\leq l$ and $1<p\leq q\leq l$, let $C_{1,n',p,q,l}^{k}$ be the $n'$-vertex $k$-uniform bicyclic hypergraph obtained from $P_{p}, P_{q}$ and $P_{l}$
by identifying three vertices $u_{1}, v_{1}, w_{1},$ and identifying three vertices $u_{p+1}, v_{q+1}, w_{l+1}$.
For $q>1, 1\leq p\leq q-1\leq l$ and $q=1, 1<p\leq l$, let $C_{2,n',p,q,l}^{k}$ be the $n'$-vertex $k$-uniform bicyclic hypergraph obtained from $P_{p}, P_{q}$ and
$P_{l}$ by identifying three vertices $u_{1}, v_{1}, w_{1}$, identifying $u_{p+1}$ with $v_{q+1}$, and identifying
$w_{l+1}$ with $v$, respectively, where $v\in f_{q}\setminus \{v_{q}, v_{q+1}\}$.
For $q>2, 1\leq p\leq q-2\leq l$ and $q=2,1\leq p\leq l$ and $q=1, k>3,1<p\leq l$, let $C_{3,n',p,q,l}^{k}$ be the $n'$-vertex $k$-uniform bicyclic hypergraph obtained from $P_{p}, P_{q}$ and
$P_{l}$ by identifying $u_{1}$ with $v_{1}$, identifying $u_{p+1}$ with $v_{q+1}$, identifying
$w_{1}$ with $v'$, and identifying $w_{l+1}$ with $v''$, respectively, where $v'\in f_{1}\setminus \{v_{1}, v_{2}\}$ and $v''\in f_{q}\setminus \{v_{q}, v_{q+1}\}$.
(in the special case $q=1$, we choose $v'\neq v''$).

For $n\geq n'$ and $i\in \{1, 2, 3\}$, let $\mathcal{C}_{i,n,p,q,l}^{k}$ be the set of $n$-vertex $m$-edge $k$-uniform bicyclic hypergraphs
each of which contains $C_{i,n',p,q,l}^{k}$ as a sub-hypergraph, where $m=\frac{n+1}{k-1}$.
%And $C_{i,n',p,q,l}^{k}$ is called the base of $\mathcal{C}_{i,n,p,q,l}^{k}$.
Let $\mathcal{C}_{n}^{k}=\{\mathcal{C}_{1,n,p,q,l}^{k} ~|~ p=1,1<q\leq l \text{~or~} 1<p\leq q\leq l\} \bigcup \{\mathcal{C}_{2,n,p,q,l}^{k}~|~ q=1,1<p\leq l \text{~or~} q>1, 1\leq p\leq q-1\leq l\} \bigcup \{\mathcal{C}_{3,n,p,q,l}^{k}~|~ q>2, 1\leq p\leq q-2\leq l\text{~or~} q=2,1\leq p\leq l  \text {~or~}q=1, k>3,1<p\leq l\}$.
Moreover, for $i\in \{1, 2\}$, let $C_{i,n,p,q,l}^{k}(m-p-q-l)$  denote the $k$-uniform bicyclic hypergraph
obtained from $C_{i,n',p,q,l}^{k}$ by attaching $m-p-q-l$ pendant edges at the vertex with
degree 3, where $m=\frac{n+1}{k-1}$.

%Let $\mathbb{B}_{n}^{k}$
%Note that $\mathcal{B}_{n}^{k}\bigcup \mathcal{C}_{n}^{k}$ is the set of all linear bicyclic $k$-uniform hypergraphs with $n$ vertices and $m$ edges.






%For $E'\subseteq E(\mathcal{H})$, let $\mathcal{H}-E'$ denote the sub-hypergraph of $\mathcal{H}$ obtained by deleting the edges of $E'$.
%Similar to transformations on graphs in \cite{2007A},
The following Lemma gives an operation of moving edges which changes the Zagreb indices.

%\begin{lem}\label{sp2}\cite{}
%Let $\mathcal{H}$ be a uniform hypergraph with $u,v\in V(\mathcal{H})$ and $d_{\mathcal{H}}(u)\geq 2$. Let $e_{1},\ldots,e_{t}$ be the pendant edges incident with $u$, $v\notin e_{i}$ for each $i\in[t]$  and $d_{\mathcal{H}}(v)>d_{\mathcal{H}}(u)-t$.  Write $e'_{i}=(e_{i}\setminus \{u\})\bigcup\{v\}$ for each $i\in[t]$. Let $\mathcal{H}^{'}$ be the hypergraph with $V(\mathcal{H}^{'})=V(\mathcal{H})$ and $E(\mathcal{H}^{'})=(E(\mathcal{H})\setminus\{e_{1},\ldots,e_{t}\})\bigcup\{e'_{1},\ldots,e'_{t}\}$.
%Then $\mathcal{H}^{'}$ is obtained from $\mathcal{H}$ by moving $t$ pendant edges from $u$ to $v$ and $M(\mathcal{H}')>M(\mathcal{H})$.
%\end{lem}

\begin{lem}\label{sp277}
For $k\geq 3$, let $\mathcal{H}$ be a linear $k$-uniform hypergraph with $u,v\in V(\mathcal{H})$ and $d_{\mathcal{H}}(u)\geq 2$. Let $e_{1},\ldots,e_{t}$ be the edges incident with $u$, $v\notin e_{i}$ for each $i\in[t]$  and $d_{\mathcal{H}}(v)>d_{\mathcal{H}}(u)-t$.  Write $e'_{i}=(e_{i}\setminus \{u\})\bigcup\{v\}$ for each $i\in[t]$. Let $\mathcal{H}^{'}$ be the hypergraph with $V(\mathcal{H}^{'})=V(\mathcal{H})$ and $E(\mathcal{H}^{'})=(E(\mathcal{H})\setminus\{e_{1},\ldots,e_{t}\})\bigcup\{e'_{1},\ldots,e'_{t}\}$.
Then $\mathcal{H}^{'}$ is obtained from $\mathcal{H}$ by moving $t$ edges from $u$ to $v$ and $M(\mathcal{H}')>M(\mathcal{H})$.
\end{lem}
\begin{proof}
By the definition of the Zagreb index, we have
\begin{align*}
M(\mathcal{H}')-M(\mathcal{H})&=d^{2}_{\mathcal{H}'}(v)+d^{2}_{\mathcal{H}'}(u)-d^{2}_{\mathcal{H}}(v)-d^{2}_{\mathcal{H}}(u)\\
&=(d_{\mathcal{H}}(v)+t)^{2}+(d_{\mathcal{H}}(u)-t)^{2}-d_{\mathcal{H}}^{2}(v)-d_{\mathcal{H}}^{2}(u)\\
&=2t(t+d_{\mathcal{H}}(v)-d_{\mathcal{H}}(u))>0.
\end{align*}
\end{proof}



%The following lemma gave the hypergraph with minimum Zagreb index among all connected $k$-uniform hypergraphs with $n$ vertices and $q$ edges.
%%Similar to graphs \cite{2004Thhe}, we get the maximum degree of the connected uniform hypergraph with minimum Zagreb index is less than or equal to $2$.
%\begin{lem}\label{ppp}\cite{2309.16925}
%In connected $k$-uniform hypergraphs with $n$ vertices and $q$ edges, the hypergraph with maximum degree $1$ or $2$ has minimum Zagreb index.
%\end{lem}




\section{Main results}

In this section,
we give the hypergraphs with the maximum and minimum Zagreb indices among all linear bicyclic uniform hypergraphs.


%By Lemma \ref{ppp}, we
The following Theorem gives the hypergraph with minimum Zagreb index among all linear bicyclic uniform hypergraphs.
\begin{thm}
In linear bicyclic $k$-uniform hypergraphs with $n$ vertices and $m$ edges, the hypergraph with maximum degree $2$ has minimum Zagreb index.
\end{thm}
\begin{proof}
%Let $\textbf{H}$ be the set of all connected $m$-uniform hypergraphs with $n$ vertices and $q$ edges. For each $\mathcal{H}\in\textbf{H}$, let $n_{k}$ be the number of vertices of $\mathcal{H}$  whose degree is equal to $k$, and $\Delta_{\mathcal{H}}$ be the maximum degree of $\mathcal{H}$.
Let $\mathcal{H}$ be a bicyclic $k$-uniform hypergraph with $n$ vertices and $m$ edges. Let $n_{t}$ be the number of vertices of $\mathcal{H}$  whose degree is equal to $t$, and $\Delta_{\mathcal{H}}$ be the maximum degree of $\mathcal{H}$. Then
\begin{equation*}\label{zsyzh1}
\sum\limits_{t=1}^{\Delta_{\mathcal{H}}}n_{t}=n,\sum\limits_{t=1}^{\Delta_{\mathcal{H}}}tn_{t}=km, \text{~and~} M(\mathcal{H})=\sum\limits_{t=1}^{\Delta_{\mathcal{H}}}t^{2}n_{t}.
\end{equation*}


By the above Equations, we have
$$
M(\mathcal{H})%=\sum\limits_{t=1}^{\Delta_{\mathcal{H}}}t^{2}n_{t}
=\sum\limits_{t=1}^{\Delta_{\mathcal{H}}}((t-1)(t-2)+3t-2)n_{t}
%&=\sum\limits_{t=1}^{\Delta_{\mathcal{H}}}(t-1)(t-2)n_{t}+3\sum\limits_{t=1}^{\Delta_{\mathcal{H}}}tn_{t}-2\sum\limits_{t=1}^{\Delta_{\mathcal{H}}}n_{t}\\
=\sum\limits_{t=1}^{\Delta_{\mathcal{H}}}(t-1)(t-2)n_{t}+3km-2n.\\
$$
Therefore, when $\Delta_{\mathcal{H}}=2$, $\mathcal{H}$ has minimum Zagreb index, and $M(\mathcal{H})=3km-2n$.

\end{proof}

The following Corollary follows immediately from the above Theorem.
\begin{cor}
In linear bicyclic $k$-uniform hypergraphs with $n$ vertices, $m$ edges and girth $g$, the hypergraph with maximum degree $2$ has minimum Zagreb index.
\end{cor}


The following Theorem
gives the hypergraph with maximum Zagreb index among all hypergraphs in $\bigcup_{i=1}^{3}\{\mathcal{B}_{i,n,g,l,q}^{k}~|~q\geq g,l\geq 0\}$, and
gives the hypergraph with maximum Zagreb index among all hypergraphs in $\mathcal{B}_{n}^{k}$.
\begin{thm}
%For $m\geq3$, in an $S$-order of $\textbf{U}_{g,f}$ the last hypergraph is $F_{g,f}$.%����ֻ��һ��Ȧ
For $k\geq3$ and $m\geq 2g$, $B_{1,n,g,0,g}^{k}(m-2g)$ is the hypergraph with maximum Zagreb index among all hypergraphs in $\bigcup_{i=1}^{3}\{\mathcal{B}_{i,n,g,l,q}^{k}~|~q\geq g,l\geq 0\}$.
And for $m\geq6$,
$B_{1,n,3,0,3}^{k}(m-6)$ is the hypergraph with maximum Zagreb index among all  hypergraphs in $\mathcal{B}_{n}^{k}$.
\end{thm}
\begin{proof}


Firstly, we consider the hypergraph  in $\mathcal{B}_{1,n,g,l,q}^{k}$. When $l>0$,
repeating the operation of moving edges of Lemma \ref{sp277}, any hypergraph in $\mathcal{B}_{1,n,g,l,q}^{k}$ can be changed into a $k$-uniform bicyclic hypergraph $\mathcal{H}_{1}$
obtained from $B_{1,n',g,l,q}^{k}$ by attaching $m-g-q-l$ pendant edges at a vertex with degree 3.
And each application of Lemma \ref{sp277} strictly increases the Zagreb index. Without loss of generality, let $d_{\mathcal{H}_{1}}(v_{2,2})=3, d_{\mathcal{H}_{1}}(v_{1,2})\geq3$.
Suppose $\mathcal{H}_{2}$ is obtained from $\mathcal{H}_{1}$ by moving 2 edges incident with $v_{2,2}$ in $E(C_{2})$
from vertex $v_{2,2}$ to vertex $v_{1,2}$. By Lemma \ref{sp277}, we have $M(\mathcal{H}_{2})>M(\mathcal{H}_{1})$.
% We have
%\begin{align*}
%  M(\mathcal{H}_{2})-M(\mathcal{H}_{1})&=d_{\mathcal{H}_{2}}^{2}(u)+d_{\mathcal{H}_{2}}^{2}(v)-d_{\mathcal{H}_{1}}^{2}(u)-d_{\mathcal{H}_{1}}^{2}(v)\\
%   &=1^{2}+(d_{\mathcal{H}_{1}}(v)+2)^{2}-3^{2}-d_{\mathcal{H}_{1}}^{2}(v)\\
%   &=4d_{\mathcal{H}_{1}}(v)-4>0.
%\end{align*}
Repeating the operation of moving edges of Lemma \ref{sp277}, $\mathcal{H}_{2}$ be changed into $B_{1,n,g,0,q}^{k}(m-g-q)$.
When $l=0$, repeating the operation of moving edges of Lemma \ref{sp277}, any hypergraph in $\mathcal{B}_{1,n,g,0,q}^{k}$ can be changed into $B_{1,n,g,0,q}^{k}(m-g-q)$.

%Therefore, $B_{1,n,g,0,q}^{k}(m-g-q)$ is last hypergraph in an $S$-order of all hypergraphs $B_{1,n,g,l,q}^{k}$.
When $q=g+s$ and $s>0$. The hypergraph
$B_{1,n,g,0,q-1}^{k}(m-g-q+1)$ can be obtained from $B_{1,n,g,0,q}^{k}(m-g-q)$ by moving 1 edges in $E(C_{2})$
from a vertex with degree $2$ adjacent to $v_{1,2}$ to $v_{1,2}$. By Lemma \ref{sp277}, we have $M(B_{1,n,g,0,q-1}^{k}(m-g-q+1))>M(B_{1,n,g,0,q}^{k}(m-g-q))$.
%We have
%\begin{align*}
%  &M(B_{1,n,g,0,q-1}^{k}(m-g-q+1))-M(B_{1,n,g,0,q}^{k}(m-g-q))\\
%  &=(d_{B_{1,n,g,0,q}^{k}(m-g-q)}(u)+1)^{2}+1^{2}-d_{B_{1,n,g,0,q}^{k}(m-g-q)}^{2}(u)-2^{2}\\
%   &=2d_{B_{1,n,g,0,q}^{k}(m-g-q)}(u)-2>0.
%\end{align*}
Similarly, we have $M(B_{1,n,g,0,q}^{k}(m-g-q))<\cdots<M(B_{1,n,g,0,q-s+1}^{k}(m-g-q+s-1))<B_{1,n,g,0,g}^{k}(m-2g)$.
Therefore, $B_{1,n,g,0,g}^{k}(m-2g)$ is the hypergraph with maximum Zagreb index in $\{\mathcal{B}_{1,n,g,l,q}^{k}~|~l\geq0,q\geq g\}$.

Secondly, we consider the hypergraph in $\mathcal{B}_{2,n,g,l,q}^{k}$.
Repeating the operation of moving edges of Lemma \ref{sp277}, any hypergraph in $\mathcal{B}_{2,n,g,l,q}^{k}$ can be changed into a $k$-uniform bicyclic hypergraph $\mathcal{H}_{3}$
obtained from $B_{2,n',g,l,q}^{k}$ by attaching $m-g-q-l$ pendant edges at the vertex with degree 3.

When $l>0$, $\mathcal{H}_{4}$ is obtained from $\mathcal{H}_{3}$ by moving $e_{l}$
from vertex $v_{2,1}$ to vertex $v_{2,2}$. By Lemma \ref{sp277}, we have $M(\mathcal{H}_{4})>M(\mathcal{H}_{3})$. Obviously, $\mathcal{H}_{4}\in \mathcal{B}_{1,n,g,l,q}^{k}.$

% \begin{align*}
%  M(\mathcal{H}_{4})-M(\mathcal{H}_{3})&=1^{2}+3^{2}-2^{2}-2^{2}>0.
%\end{align*}
When $l=0$, $\mathcal{H}_{4}$ is obtained from $\mathcal{H}_{3}$ by moving $m-g-q+2$ edges in $E(\mathcal{H}_{3})-E(C_{2})$
from vertex $v_{2,1}$ to vertex $v_{2,2}$. By Lemma \ref{sp277}, we have $M(\mathcal{H}_{4})>M(\mathcal{H}_{3})$. Obviously, $\mathcal{H}_{4}\in \mathcal{B}_{1,n,g,0,q}^{k}.$
%\begin{align*}
%  M(\mathcal{H}_{4})-M(\mathcal{H}_{3})&=(m-g-q+4)^{2}+1^{2}-(m-g-q+3)^{2}-2^{2}\\
%  &=15+2(m-g-q)>0.
%\end{align*}
Therefore, $B_{1,n,g,0,g}^{k}(m-2g)$ is the hypergraph with maximum Zagreb index in $\bigcup_{i=1}^{2}\{\mathcal{B}_{i,n,g,l,q}^{k}~|~l\geq0,q\geq g\}$.

Thirdly, we consider the hypergraph in $\mathcal{B}_{3,n,g,l,q}^{k}$. Repeating the operation of moving edges of Lemma \ref{sp277}, any hypergraph in $\mathcal{B}_{3,n,g,l,q}^{k}$ can be changed into a $k$-uniform bicyclic hypergraph $\mathcal{H}_{5}$ obtained from $B_{3,n',g,l,q}^{k}$ by attaching $m-g-q-l$ pendant edges at the vertex with degree 2.

When $l>0$. If $\mathcal{H}_{5}$ is obtained from $B_{3,n',g,l,q}^{k}$ by attaching $m-g-q-l$ pendant edges at the vertex $v_{1,1}$. Let $\mathcal{H}_{6}$
be obtained from $\mathcal{H}_{5}$ by moving $m-g-q-l+1$ edges in $E(\mathcal{H}_{5})-E(C_{1})$
from vertex $v_{1,1}$ to vertex $v_{1,2}$. By Lemma \ref{sp277}, we have $M(\mathcal{H}_{6})>M(\mathcal{H}_{5})$. Obviously, $\mathcal{H}_{6}\in \mathcal{B}_{2,n,g,l,q}^{k}.$
If $\mathcal{H}_{5}$ is obtained from $B_{3,n',g,l,q}^{k}$ by attaching $m-g-q-l$ pendant edges at the vertex except $v_{1,1}$ with degree 2. Let $\mathcal{H}_{6}$
be obtained from $\mathcal{H}_{5}$ by moving $e_{1}$
from vertex $v_{1,1}$ to vertex $v_{1,2}$. By Lemma \ref{sp277}, we have $M(\mathcal{H}_{6})>M(\mathcal{H}_{5})$. Obviously, $\mathcal{H}_{6}\in \mathcal{B}_{2,n,g,l,q}^{k}.$

%
%be obtained by moving the vertex $u_{0}$ of $\mathcal{H}_{5}$
%from vertex $v_{1,1}$ to vertex $v_{1,2}$.
%\begin{align*}
%  M(\mathcal{H}_{6})-M(\mathcal{H}_{5})&=3^{2}+1^{2}-2^{2}-2^{2}>0.
%\end{align*}
%Let $\mathcal{H}_{7}$ be obtained by moving the vertex $u_{l}$ of $\mathcal{H}_{6}$
%from vertex $v_{2,1}$ to vertex $v_{2,2}$. Similarly, we have $M(\mathcal{H}_{7})-M(\mathcal{H}_{6})>0.$ Obviously, $\mathcal{H}_{7}\in B_{1,n,g,l,q}^{k}.$


When $l=0$.
If $\mathcal{H}_{5}$ is obtained from $B_{3,n',g,0,q}^{k}$ by attaching $m-g-q$ pendant edges at the vertex $v_{1,1}$. Let $\mathcal{H}_{7}$ be obtained from $\mathcal{H}_{5}$ by moving $m-g-q+1$ edges in $E(\mathcal{H}_{5})-E(C_{1})$
from vertex $v_{1,1}$ to vertex $v_{1,2}$. By Lemma \ref{sp277}, we have $M(\mathcal{H}_{7})>M(\mathcal{H}_{5})$. Obviously, $\mathcal{H}_{7}\in \mathcal{B}_{2,n,g,0,q}^{k}.$
If $\mathcal{H}_{5}$ is obtained from $B_{3,n',g,0,q}^{k}$ by attaching $m-g-q$ pendant edges at the vertex except $v_{1,1}$ with degree 2. Let $\mathcal{H}_{7}$ be obtained from $\mathcal{H}_{5}$ by moving $1$ edge in $E(C_{2})$
from vertex $v_{1,1}$ to vertex $v_{1,2}$. By Lemma \ref{sp277}, we have $M(\mathcal{H}_{7})>M(\mathcal{H}_{5})$. Obviously, $\mathcal{H}_{7}\in \mathcal{B}_{2,n,g,0,q}^{k}.$





%$\mathcal{H}_{8}$ be obtained by moving the vertex $u_{0}$ of $\mathcal{H}_{5}$
%from vertex $v_{1,1}$ to vertex $v_{1,2}$ and moving the vertex $u_{l}$ of $\mathcal{H}_{5}$
%from vertex $v_{2,1}$ to vertex $v_{2,2}$. Obviously, $\mathcal{H}_{8}\in B_{1,n,g,l,q}^{k}.$
%\begin{align*}
%  M(\mathcal{H}_{8})-M(\mathcal{H}_{5})&=4^{2}+1^{2}+1^{2}-2^{2}-2^{2}-2^{2}>0.
%\end{align*}
Therefore, $B_{1,n,g,0,g}^{k}(m-2g)$ is the hypergraph with maximum Zagreb index in
$\bigcup_{i=1}^{3}\{\mathcal{B}_{i,n,g,l,q}^{k}~|~l\geq0,q\geq g\}$.
%$\{B_{n,g,l,q}^{k},l\geq0,q\geq g\}$.

For $3\leq g\leq\frac{m}{2}$, we have
\begin{align*}
M(B_{1,n,g,0,g}^{k}(m-2g))&=2g(k-2)+(m-2g)(k-1)+8(g-1)+(m-2g+4)^{2}\\
&=-10g+mk+7m+8+m^{2}+4g^{2}-4mg.
\end{align*}
%$M(B_{1,n,g,0,g}^{k}(m-2g))=-10g+mk+7m+8+m^{2}+4g^{2}-4mg,3\leq g\leq\frac{m}{2}.$
Let $f(x)=-10x+mk+7m+8+m^{2}+4x^{2}-4mx, x\in[3,\frac{m}{2}].$ Since $\frac{df(x)}{dx}=-10+8x-4m<0$, $f(x)$ is a strictly monotone decreasing function. So $M(B_{1,n,g,0,g}^{k}(m-2g)\leq M(B_{1,n,3,0,3}^{k}(m-6)$ for $3\leq g\leq\frac{m}{2}$ with the equation if and only if $g=3$. Hence, $B_{1,n,3,0,3}^{k}(m-6)$ is the hypergraph with maximum Zagreb index among all hypergraphs in $\mathcal{B}_{n}^{k}=\bigcup_{i=1}^{3}\{\mathcal{B}_{i,n,p,l,q}^{k}~|~q\geq p\geq3,l\geq 0\}$.

\end{proof}

The following Theorem gives the hypergraph with maximum Zagreb index among  all hypergraphs with girth $g$ in $\mathcal{C}_{n}^{k}$, and gives the hypergraph with maximum Zagreb index among all hypergraphs in $\mathcal{C}_{n}^{k}$.
\begin{thm}\label{q1}
%For $m\geq3$, in an $S$-order of $\textbf{U}_{g,f}$ the last hypergraph is $F_{g,f}$.%����ֻ��һ��Ȧ
For $k\geq3$ and $m\geq \frac{3g}{2}$, when $g$ is even, $C_{1,n,\frac{g}{2},\frac{g}{2},\frac{g}{2}}^{k}(m-\frac{3g}{2})$ is the hypergraph with maximum Zagreb index among all hypergraphs with girth $g$ in $\mathcal{C}_{n}^{k}$.
When $g$ is odd, $C_{2,n,\lfloor\frac{g}{2}\rfloor,\lceil\frac{g}{2}\rceil,\lfloor\frac{g}{2}\rfloor}^{k}(m-g-\lfloor\frac{g}{2}\rfloor)$ is the hypergraph with maximum Zagreb index among all  hypergraphs with girth $g$ in $\mathcal{C}_{n}^{k}$.

For $m\geq 6$, $C_{2,n,1,2,1}^{k}(m-4)$ is the hypergraph with maximum Zagreb index among all hypergraphs in $\mathcal{C}_{n}^{k}$.
%where $m$ is the number of edges.
\end{thm}
\begin{proof}
Firstly, we consider the hypergraph in $\mathcal{C}_{1,n,p,g-p,l}^{k}$.
Repeating the operation of moving edges of Lemma \ref{sp277}, any hypergraph in $\mathcal{C}_{1,n,p,g-p,l}^{k}$ can be changed into a $k$-uniform bicyclic hypergraph $C_{1,n,p,g-p,l}^{k}(m-g-l)$
obtained from $C_{1,n',p,g-p,l}^{k}$ by attaching $m-g-l$ pendant edges at a vertex with degree 3.
And each application of Lemma \ref{sp277} strictly increases the Zagreb index.
%When $p>1$ and $p+q\neq l+q$. Without loss of generality, suppose $p+q\leq l+q$.
%Without loss of generality, suppose $p\leq q\leq l$, obviously, $p+q=g$.
%Let $C_{1,n,p,g-p,l}^{k}(m-g-l)$ denote the $k$-uniform bicyclic hypergraph
%obtained from $C_{1,n',p,g-p,l}^{k}$  by attaching $m-g-l$ pendant edges at the vertex with
%degree 3.

If $g$ is even. When $g-p<l$, without loss of generality, let $d_{C_{1,n,p,g-p,l}^{k}(m-g-l)}(u_{1})\geq3$.
The hypergraph
$C_{1,n,p,g-p,l-1}^{k}(m-g-l+1)$ can
be obtained from $C_{1,n,p,g-p,l}^{k}(m-g-l)$ by moving $g_{2}$
from vertex $w_{2}$ to vertex $u_{1}$. By Lemma \ref{sp277}, we have $M(C_{1,n,p,g-p,l-1}^{k}(m-g-l+1))>M(C_{1,n,p,g-p,l}^{k}(m-g-l))$.
%we have
%\begin{align*}
%  &M(C_{1,n,p,g-p,l-1}^{k}(m-g-l+1))-M(C_{1,n,p,g-p,l}^{k}(m-g-l))\\
%  &=(m-g-l+1+3)^{2}+1^{2}-(m-g-l+3)^{2}-2^{2}\\
%   &=2(m-g-l)+4>0.
%\end{align*}
Similarly, we get $M(C_{1,n,p,g-p,l}^{k}(m-g-l))<M(C_{1,n,p,g-p,l-1}^{k}(m-g-l+1))<\cdots<M(C_{1,n,p,g-p,g-p}^{k}(m-2g+p)).$

When $p\neq \frac{g}{2}$, we have
\begin{align*}
  &M(C_{1,n,p+1,g-p-1,g-p-1}^{k}(m-2g+p+1))-M(C_{1,n,p,g-p,g-p}^{k}(m-2g+p))\\
  &=(m-2g+p+4)^{2}+1^{2}-(m-2g+p+3)^{2}-2^{2}\\
   &=2(m-2g+p)+4>0.
\end{align*}
Similarly, we get $M(C_{1,n,p,g-p,g-p}^{k}(m-2g+p))<M(C_{1,n,p+1,g-p-1,g-p-1}^{k}(m-2g+p+1))<\cdots<M(C_{1,n,\frac{g}{2},\frac{g}{2},\frac{g}{2}}^{k}(m-\frac{3g}{2}))$. Therefore, when $g$ is even, $C_{1,n,\frac{g}{2},\frac{g}{2},\frac{g}{2}}^{k}(m-\frac{3g}{2})$ has maximum Zagreb index among all hypergraphs in $\{\mathcal{C}_{1,n,p,g-p,l}^{k}~|~ p=1,1<g-p\leq l \text{~or~} 1<p\leq g-p\leq l\}$.

%When $p\neq q=l$, $p=q\neq l$ or $p=q\neq l$
If $g$ is odd, similarly, we get $C_{1,n,\lfloor\frac{g}{2}\rfloor,\lceil\frac{g}{2}\rceil,\lceil\frac{g}{2}\rceil}^{k}(m-g-\lceil\frac{g}{2}\rceil)$ has maximum Zagreb index among all hypergraphs  in $\{\mathcal{C}_{1,n,p,g-p,l}^{k}~|~ p=1,1<g-p\leq l \text{~or~} 1<p\leq g-p\leq l\}$.

Secondly, for $p+q=g$, we consider the hypergraph in $\mathcal{C}_{2,n,p,q,l}^{k}$. Repeating the operation of moving edges of Lemma \ref{sp277}, any hypergraph in $\mathcal{C}_{2,n,p,q,l}^{k}$ can be changed into a $k$-uniform bicyclic hypergraph $\mathcal{H}_{1}$
obtained from $C_{2,n',p,q,l}^{k}$ by attaching $m-p-q-l$ pendant edges at the vertex with degree 3.

If $q$ of $\mathcal{H}_{1}$ is equal to $1$, then the girth is $p+1$.
$\mathcal{H}_{2}$ is obtained from $\mathcal{H}_{1}$ by moving $g_{l}$
from vertex $v$ to vertex $v_{2}$. Obviously, $\mathcal{H}_{2}\in \mathcal{C}_{1,n,1,p,l}^{k}$ and $g(\mathcal{H}_{2})=p+1$. By Lemma \ref{sp277}, we have $ M(\mathcal{H}_{2})>M(\mathcal{H}_{1})$.
%\begin{align*}
%  M(\mathcal{H}_{2})-M(\mathcal{H}_{1})&=3^{2}+1^{2}-2^{2}-2^{2}>0.
%\end{align*}

If $q$ of $\mathcal{H}_{1}$ is greater than 1, then $1\leq p\leq q-1\leq l$.

When $1\leq p\leq q-1=l$.
If $1\leq p<q-1=l$, $\mathcal{H}_{2}'$ is obtained from $\mathcal{H}_{1}$ by moving $g_{l}$
from vertex $v$ to vertex $v_{q}$. Obviously, $\mathcal{H}_{2}'\in \mathcal{C}_{1,n,p+1,q-1,l}^{k}$ and $g(\mathcal{H}_{2}')=p+q$. By Lemma \ref{sp277}, we have $ M(\mathcal{H}_{2}')>M(\mathcal{H}_{1})$.
%\begin{align*}
%  M(\mathcal{H}_{2})-M(\mathcal{H}_{1})&=3^{2}+1^{2}-2^{2}-2^{2}>0.
%\end{align*}

If $1\leq p=q-1=l$, then $g(\mathcal{H}_{1})=2l+1$. Obviously, the girth of $\mathcal{H}_{1}$ is odd. And $\mathcal{H}_{1}=C_{2,n,l,l+1,l}^{k}(m-3l-1)=C_{2,n,\lfloor\frac{g}{2}\rfloor,\lceil\frac{g}{2}\rceil,\lfloor\frac{g}{2}\rfloor}^{k}(m-g-\lfloor\frac{g}{2}\rfloor)$. We have
\begin{align*}
&M(C_{2,n,\lfloor\frac{g}{2}\rfloor,\lceil\frac{g}{2}\rceil,\lfloor\frac{g}{2}\rfloor}^{k}(m-g-\lfloor\frac{g}{2}\rfloor))\\
&=(g+\frac{g-1}{2}-1)(k-2)+(k-3)+(m-g-\frac{g-1}{2})(k-1)+4(g+\frac{g-1}{2}-1)\\
&+(3+m-g-\frac{g-1}{2})^{2}\\
&=-6g+\frac{23}{4}+mk+6m+m^{2}+\frac{9g^{2}}{4}-3mg.
\end{align*}
%\begin{align*}
%M(\mathcal{H}_{1})&=3l(k-2)+k-3+(k-1)(m-3l-1)+12l+(m-3l+2)^{2}\\
%&=-3l+2+mk+3m+m^{2}+9l^{2}-6ml.
%\end{align*}
%$$
%M(\mathcal{H}_{1})=-3l+2+mk+3m+m^{2}+9l^{2}-6ml.
%$$

When the girth is odd, $C_{1,n,\lfloor\frac{g}{2}\rfloor,\lceil\frac{g}{2}\rceil,\lceil\frac{g}{2}\rceil}^{k}(m-g-\lceil\frac{g}{2}\rceil)$ has maximum Zagreb index among all  hypergraphs  in $\{\mathcal{C}_{1,n,p,g-p,l}^{k}~|~ p=1,1<g-p\leq l \text{~or~} 1<p\leq g-p\leq l\}$. We have
\begin{align*}
&M(C_{1,n,\lfloor\frac{g}{2}\rfloor,\lceil\frac{g}{2}\rceil,\lceil\frac{g}{2}\rceil}^{k}(m-g-\lceil\frac{g}{2}\rceil))\\
&=(g+\frac{g+1}{2})(k-2)+(m-g-\frac{g+1}{2})(k-1)+4(g+\frac{g+1}{2}-3)+9+(3+m-g-\frac{g+1}{2})^{2}\\
&=-3g+\frac{19}{4}+mk+4m+m^{2}+\frac{9}{4}g^{2}-3mg.
\end{align*}
%Hence,
%$$
%M(C_{1,n,\lfloor\frac{2l+1}{2}\rfloor,\lceil\frac{2l+1}{2}\rceil,\lceil\frac{2l+1}{2}\rceil}^{k}(m-(2l+1)-\lceil\frac{2l+1}{2}\rceil))=4+mk+m+m^{2}+9l^{2}+3l-6ml.
%$$
%When $m>3l+1$, $M(C_{1,n,\lfloor\frac{2l+1}{2}\rfloor,\lceil\frac{2l+1}{2}\rceil,\lceil\frac{2l+1}{2}\rceil}^{k}(m-(2l+1)-\lceil\frac{2l+1}{2}\rceil))-M(\mathcal{H}_{1})=2-2m+6l<0$.
%When $m>\frac{3}{2}g-\frac{1}{2}$,
Hence, $M(C_{1,n,\lfloor\frac{g}{2}\rfloor,\lceil\frac{g}{2}\rceil,\lceil\frac{g}{2}\rceil}^{k}(m-g-\lceil\frac{g}{2}\rceil))-M(C_{2,n,\lfloor\frac{g}{2}\rfloor,\lceil\frac{g}{2}\rceil,\lfloor\frac{g}{2}\rfloor}^{k}(m-g-\lfloor\frac{g}{2}\rfloor))=3g-1-2m<0$.

When $1\leq p\leq q-1<l$,
$\mathcal{H}_{2}''$ is obtained from $\mathcal{H}_{1}$ by moving $g_{l}$
from vertex $v$ to vertex $v_{q+1}$. Obviously, $\mathcal{H}_{2}''\in \mathcal{C}_{1,n,p,g-p,l}^{k}$ and $g(\mathcal{H}_{2}'')=p+q$. By Lemma \ref{sp277}, we have $ M(\mathcal{H}_{2}'')>M(\mathcal{H}_{1})$.

%When $p+q>p+l+1$,
%$\mathcal{H}_{2}$ is obtained from $\mathcal{H}_{1}$ by moving $g_{l}$
%from vertex $v$ to vertex $v_{q}$. Obviously, $\mathcal{H}_{2}\in C_{1,n,p+1,l,q-1}^{k}$ or $C_{1,n,l,p+1,q-1}^{k}$ and $g(\mathcal{H}_{2})=p+l+1$. By Lemma \ref{sp277}, we have $ M(\mathcal{H}_{2})>M(\mathcal{H}_{1})$.


Therefore,
when $g$ is even, $C_{1,n,\frac{g}{2},\frac{g}{2},\frac{g}{2}}^{k}(m-\frac{3g}{2})$ is the hypergraph with maximum Zagreb index among all hypergraphs with girth $g$ in $\{\mathcal{C}_{1,n,p,q,l}^{k} ~|~ p=1,1<q\leq l \text{~or~} 1<p\leq q\leq l\} \bigcup \{\mathcal{C}_{2,n,p,q,l}^{k}~|~ q=1,1<p\leq l \text{~or~} q>1, 1\leq p\leq q-1\leq l\}$. When $g$ is odd, $C_{2,n,\lfloor\frac{g}{2}\rfloor,\lceil\frac{g}{2}\rceil,\lfloor\frac{g}{2}\rfloor}^{k}(m-g-\lfloor\frac{g}{2}\rfloor)$ is the hypergraph with maximum Zagreb index among all hypergraphs with girth $g$ in $\{\mathcal{C}_{1,n,p,q,l}^{k} ~|~ p=1,1<q\leq l \text{~or~} 1<p\leq q\leq l\} \bigcup \{\mathcal{C}_{2,n,p,q,l}^{k}~|~ q=1,1<p\leq l \text{~or~} q>1, 1\leq p\leq q-1\leq l\} $.

Thirdly, for $p+q=g$, we consider the hypergraph in $\mathcal{C}_{3,n,p,q,l}^{k}$. Repeating the operation of moving edges of Lemma \ref{sp277}, any hypergraph in $\mathcal{C}_{3,n,p,q,l}^{k}$ can be changed into a $k$-uniform bicyclic hypergraph $\mathcal{H}_{3}$
obtained from $C_{3,n',p,q,l}^{k}$ by attaching $m-p-q-l$ pendant edges at a vertex with degree 2.


If $q$ of $\mathcal{H}_{3}$ is equal to 1, then $g(\mathcal{H}_{3})=p+1$.
If $\mathcal{H}_{3}$ is obtained from $C_{3,n',p,1,l}^{k}$ by attaching $m-p-1-l$ pendant edges at the vertex $v'$ (or $v''$). Let $\mathcal{H}_{4}$ be obtained from $\mathcal{H}_{3}$ by moving $m-p-1-l$ pendant edges from vertex $v'$ to vertex $v_{1}$ (or from vertex $v''$ to vertex $v_{2}$), moving $g_{1}$
from vertex $v'$ to vertex $v_{1}$ and moving $g_{l}$
from vertex $v''$ to vertex $v_{2}$. Obviously, $\mathcal{H}_{4}\in \mathcal{C}_{1,n,1,p,l}^{k}$ and $g(\mathcal{H}_{4})=p+1$.
By Lemma \ref{sp277}, we have $M(\mathcal{H}_{4})>M(\mathcal{H}_{3})$.
If $\mathcal{H}_{3}$ is obtained from $C_{3,n',p,1,l}^{k}$ by attaching $m-p-1-l$ pendant edges at the vertex except $v',v''$ with degree 2.  Let $\mathcal{H}_{4}$ be obtained from $\mathcal{H}_{3}$ by moving $g_{1}$ from vertex $v'$ to vertex $v_{1}$ and moving $g_{l}$
from vertex $v''$ to vertex $v_{2}$. Obviously, $\mathcal{H}_{4}\in \mathcal{C}_{1,n,1,p,l}^{k}$ and $g(\mathcal{H}_{4})=p+1$. By Lemma \ref{sp277}, we have $M(\mathcal{H}_{4})>M(\mathcal{H}_{3})$.
%\begin{align*}
%  M(\mathcal{H}_{4})-M(\mathcal{H}_{3})&=3^{2}+3^{2}+1^{2}+1^{2}-2^{2}-2^{2}-2^{2}-2^{2}>0.
%\end{align*}


If $q=2, p=l=1$ of $\mathcal{H}_{3}$, then $g(\mathcal{H}_{3})=3$.
If $\mathcal{H}_{3}$ is obtained from $C_{3,n',1,2,1}^{k}$ by attaching $m-4$ pendant edges at the vertex $v'$. Let $\mathcal{H}_{5}$ be obtained from $\mathcal{H}_{3}$ by moving $g_{1}$ and $m-4$ pendant edges from vertex $v'$ to vertex $v_{1}$. Obviously, $\mathcal{H}_{5}\in \mathcal{C}_{2,n,1,2,1}^{k}$ and $g(\mathcal{H}_{5})=3$. By Lemma \ref{sp277}, we have $M(\mathcal{H}_{5})>M(\mathcal{H}_{3})$.
If $\mathcal{H}_{3}$ is obtained from $C_{3,n',1,2,1}^{k}$ by attaching $m-4$ pendant edges at the vertex except $v'$ with degree 2.  Let $\mathcal{H}_{5}$ be obtained from $\mathcal{H}_{3}$ by moving $g_{1}$
from vertex $v'$ to vertex $v_{1}$. Obviously, $\mathcal{H}_{5}\in \mathcal{C}_{2,n,1,2,1}^{k}$ and $g(\mathcal{H}_{5})=3$. By Lemma \ref{sp277}, we have $M(\mathcal{H}_{5})>M(\mathcal{H}_{3})$.


If $q=2, 1=p<l$ (or $q=2, 1<p\leq l$) of $\mathcal{H}_{3}$, then $g(\mathcal{H}_{3})=3$ (or $p+2$). Similar to the proof of $q=1$ of $\mathcal{H}_{3}$, $\mathcal{H}_{3}$ can be changed into $\mathcal{H}_{4}$,
$\mathcal{H}_{4}\in \mathcal{C}_{1,n,1,2,l}^{k}$ (or $\mathcal{C}_{1,n,2,p,l}^{k}$), $g(\mathcal{H}_{4})=3$ (or $p+2$) and $M(\mathcal{H}_{4})>M(\mathcal{H}_{3})$.



If $q>2$ of $\mathcal{H}_{3}$, then $1\leq p\leq q-2\leq l$.
When $q\geq 2$, $M(\mathcal{H}_{3})=-p-q+2-l+mk+3m+m^{2}+p^{2}+q^{2}+l^{2}-2mp-2mq-2ml+2pq+2pl+2ql.$

When $1\leq p\leq q-2=l$, if $1\leq p=q-2=l$, then $M(\mathcal{H}_{3})=9l+mk-m+9l^{2}-6ml+m^{2}+4$ and $g(\mathcal{H}_{3})=2l+2.$

When $g$ is even, $C_{1,n,\frac{g}{2},\frac{g}{2},\frac{g}{2}}^{k}(m-\frac{3g}{2})$ is the hypergraph with maximum Zagreb index among all hypergraphs with girth $g$ in $\{\mathcal{C}_{1,n,p,q,l}^{k} ~|~ p=1,1<q\leq l \text{~or~} 1<p\leq q\leq l\} \bigcup \{\mathcal{C}_{2,n,p,q,l}^{k}~|~ q=1,1<p\leq l \text{~or~} q>1, 1\leq p\leq q-1\leq l\} $.
We have
\begin{align*}
M(C_{1,n,\frac{g}{2},\frac{g}{2},\frac{g}{2}}^{k}(m-\frac{3g}{2}))&=\frac{3}{2}g(k-2)+(m-\frac{3}{2}g)(k-1)+12(\frac{g}{2}-1)+9+(3+m-\frac{3}{2}g)^{2}\\
&=-\frac{9}{2}g+mk+5m+6+m^{2}+\frac{9}{4}g^{2}-3mg.
\end{align*}
Hence, $M(C_{1,n,l+1,l+1,l+1}^{k}(m-3l-3))=9l+mk-m+6+m^{2}+9l^{2}-6ml.$ Therefore, $M(C_{1,n,l+1,l+1,l+1}^{k}(m-3l-3))-M(\mathcal{H}_{3})=2>0$.

If $1\leq p<q-2=l$,
%If $\mathcal{H}_{3}$ is obtained from $C_{3,n',p,q,l}^{k}$ by attaching $m-p-q-l$ pendant edges at the vertex $v_{2}$. Let $\mathcal{H}_{4}$ be obtained by moving the vertex $w_{1}$ of $\mathcal{H}_{3}$
%from vertex $v'$ to vertex $v_{2}$. Obviously, $\mathcal{H}_{4}\in C_{2,n,p+1,q-1,l}^{k}$ and $g(\mathcal{H}_{4})=p+q$.
%\begin{align*}
%  M(\mathcal{H}_{4})-M(\mathcal{H}_{3})&=(d_{\mathcal{H}_{3}}(v_{2})+1)^{2}+1^{2}-d_{\mathcal{H}_{3}}^{2}(v_{2})-2^{2}\\
%  &=2d_{\mathcal{H}_{3}}(v_{2})-2>0.
%\end{align*}
supppose $\mathcal{H}_{3}$ is obtained from $C_{3,n',p,q,l}^{k}$ by attaching $m-p-q-l$ pendant edges at the vertex $v'$. Let $\mathcal{H}_{6}$ be obtained from $\mathcal{H}_{3}$ by moving $g_{1}$ and $m-p-q-l$ pendant edges from vertex $v'$ to vertex $v_{2}$. Obviously, $\mathcal{H}_{6}\in \mathcal{C}_{2,n,p+1,q-1,l}^{k}$ and $g(\mathcal{H}_{6})=p+q$. By Lemma \ref{sp277}, we have $M(\mathcal{H}_{6})>M(\mathcal{H}_{3})$.
%\begin{align*}
%  M(\mathcal{H}_{4})-M(\mathcal{H}_{3})&=(2+d_{\mathcal{H}_{3}}(v')-1)^{2}+1^{2}-d_{\mathcal{H}_{3}}^{2}(v')-2^{2}\\
%  &=2d_{\mathcal{H}_{3}}(v')-2>0.
%\end{align*}
If $\mathcal{H}_{3}$ is obtained from $C_{3,n',p,q,l}^{k}$ by attaching $m-p-q-l$ pendant edges at the vertex except $v'$ with degree 2.  Let $\mathcal{H}_{6}$ be obtained from $\mathcal{H}_{3}$ by moving $g_{1}$
from vertex $v'$ to vertex $v_{2}$. Obviously, $\mathcal{H}_{6}\in \mathcal{C}_{2,n,p+1,q-1,l}^{k}$ and $g(\mathcal{H}_{6})=p+q$. By Lemma \ref{sp277}, we have $M(\mathcal{H}_{6})>M(\mathcal{H}_{3})$.
%\begin{align*}
%  M(\mathcal{H}_{4})-M(\mathcal{H}_{3})&=3^{2}+1^{2}-2^{2}-2^{2}>0.
%\end{align*}

When $1\leq p\leq q-2<l$, similar to the proof of $q=2, p=l=1$ of $\mathcal{H}_{3}$, $\mathcal{H}_{3}$ can be changed into $\mathcal{H}_{5}$,
$\mathcal{H}_{5}\in \mathcal{C}_{2,n,p,q,l}^{k}$, $g(\mathcal{H}_{5})=p+q$ and $M(\mathcal{H}_{5})>M(\mathcal{H}_{3})$.


%When $p+q>p+l+2$, similar to the proof of $q>2, p+q=p+l+2, p<l$ of $\mathcal{H}_{3}$, $\mathcal{H}_{3}$ can be changed into $\mathcal{H}_{4}$, $M(\mathcal{H}_{4})>M(\mathcal{H}_{3})$,
%$\mathcal{H}_{4}\in C_{2,n,p+1,q-1,l}^{k}$ or $C_{2,n,l,q-1,p+1}^{k}$ and $g(\mathcal{H}_{4})=p+l+2$.




%$C_{1,n,p,g-p,l}^{k}\bigcup C_{2,n,p,q,l}^{k}\bigcup C_{3,n,p,q,l}^{k}$.


Therefore,
when $g$ is even, $C_{1,n,\frac{g}{2},\frac{g}{2},\frac{g}{2}}^{k}(m-\frac{3g}{2})$ is the hypergraph with maximum Zagreb index among all hypergraphs with girth $g$ in $\mathcal{C}_{n}^{k}$.
When $g$ is odd, $C_{2,n,\lfloor\frac{g}{2}\rfloor,\lceil\frac{g}{2}\rceil,\lfloor\frac{g}{2}\rfloor}^{k}(m-g-\lfloor\frac{g}{2}\rfloor)$ is the hypergraph with maximum Zagreb index among all hypergraphs with girth $g$ in $\mathcal{C}_{n}^{k}$.



If $\mathcal{H} \in \{\mathcal{C}_{1,n,p,q,l}^{k} ~|~ p=1,1<q\leq l \text{~or~} 1<p\leq q\leq l\}$.
Since $g=p+q$ and $p \leq q$, $g-q \leq q$, that is $q\geq \frac{g}{2}$. Since $l\geq q$, we have $l\geq \frac{g}{2}$.
Then $p+q+l=g+l\geq \frac{3}{2}g$. Thus,
when $m\geq \frac{3}{2}g$, $\mathcal{H}$ exists.

If $\mathcal{H} \in  \{\mathcal{C}_{2,n,p,q,l}^{k}~|~ q=1,1<p\leq l \text{~or~} q>1, 1\leq p\leq q-1\leq l\}$.
Since $g=p+q$,
we have $p+q\leq p+l+1$ and $p+q\leq q+l$,
which implies $l\geq \frac{g-1}{2}$.
Then we have $p+q+l\geq g+\frac{g-1}{2}=\frac{3}{2}g-\frac{1}{2}$.
Thus, when
$m\geq \frac{3}{2}g-\frac{1}{2}$, $\mathcal{H}$ exists.
For $\frac{3g}{2}-\frac{1}{2}\leq m<\frac{3g}{2}$, we have $m=\frac{3g}{2}-\frac{1}{2}$ and $g$ is odd.
If $q=1$, $\mathcal{H}$ does not exist.
If $q>1$, $1\leq g-q\leq q-1\leq l$, then $\frac{g}{2}+\frac{1}{2}\leq q\leq l+1.$  Since $m=\frac{3g}{2}-\frac{1}{2}$, $l\geq \frac{g-1}{2}$ and $\frac{g}{2}+\frac{1}{2}\leq q\leq l+1$, $l=\frac{g-1}{2}$ and $q=\frac{g}{2}+\frac{1}{2}$.
Therefore, $\mathcal{H}=C_{2,n,\lfloor\frac{g}{2}\rfloor,\lceil\frac{g}{2}\rceil,\lfloor\frac{g}{2}\rfloor}^{k}$ and $g$ is odd. %We have $M(\mathcal{H})=m(k-2)-1+4(m-1)+9=2m+4+mk$.
%The hypergraph $\mathcal{H}$ contains $m(k-2)-1$ vertices of degree $1$, $m-1$ vertices of degree $2$, and exactly one vertex of degree $3$.
%Hence, the Zagreb indices of all hypergraphs in $\{C_{2,n,g-q,q,\frac{g}{2}-\frac{1}{2}}^{k}~|~q\geq \frac{g}{2}+\frac{1}{2}, ~g~ \text{is odd}\}$ are equal.
%We have $C_{2,n,\lfloor\frac{g}{2}\rfloor,\lceil\frac{g}{2}\rceil,\lfloor\frac{g}{2}\rfloor}^{k}\in \{C_{2,n,g-q,q,\frac{g}{2}-\frac{1}{2}}^{k}~|~q\geq \frac{g}{2}+\frac{1}{2}, ~g~ \text{is odd}\}$.
%We have $C_{2,n,\frac{g}{2}-\frac{1}{2},\frac{g}{2}+\frac{1}{2},\frac{g}{2}-\frac{1}{2}}^{k}\in \{C_{2,n,g-q,q,\frac{g}{2}-\frac{1}{2}}^{k}~|~q\geq \frac{g}{2}+\frac{1}{2}, ~g~ \text{is odd}\}$.

If $\mathcal{H} \in \{\mathcal{C}_{3,n,p,q,l}^{k}~|~ q>2, 1\leq p\leq q-2\leq l\text{~or~} q=2,1\leq p\leq l  \text {~or~}q=1, k>3,1<p\leq l\}$.
Since $g=p+q$,
we have $p+q\leq p+l+2$ and $p+q\leq q+l$,
which implies $l\geq \frac{1}{2}g-1$.
Then we have $p+q+l\geq g+\frac{1}{2}g-1=\frac{3}{2}g-1$.
Thus,
when $m\geq \frac{3}{2}g-1$, $\mathcal{H}$ exists.
If $q>2$, $1\leq g-q\leq q-2\leq l$, then $\frac{g}{2}+1\leq q \leq l+2.$
For $\frac{3g}{2}-1\leq m<\frac{3g}{2}$, when $g$ is even, $m=\frac{3g}{2}-1$.
Since $m=\frac{3g}{2}-1$, $l\geq \frac{1}{2}g-1$ and $\frac{g}{2}+1\leq q \leq l+2$, $l=\frac{1}{2}g-1$ and $q=\frac{g}{2}+1$.
Therefore, $\mathcal{H}=C_{3,n,\frac{g}{2}-1,\frac{g}{2}+1,\frac{g}{2}-1}^{k}$ and $g$ is even.
When $g$ is odd, $m=\frac{3g}{2}-\frac{1}{2}$. Since $m=\frac{3g}{2}-\frac{1}{2}$, $l\geq \frac{1}{2}g-1$ and $\frac{g}{2}+1\leq q \leq l+2$, $l=\frac{1}{2}g-\frac{1}{2}$ and $q=\frac{1}{2}g+\frac{3}{2}$. Therefore, $\mathcal{H}=C_{3,n,\frac{1}{2}g-\frac{3}{2},\frac{1}{2}g+\frac{3}{2},\frac{1}{2}g-\frac{1}{2}}^{k}$ and $g$ is odd.
If $q=1$ or $2$, for $\frac{3g}{2}-1\leq m<\frac{3g}{2}$ and $m\geq 6$, $\mathcal{H}$ does not exist.





When $g$ is even, we have
$M(C_{1,n,\frac{g}{2},\frac{g}{2},\frac{g}{2}}^{k}(m-\frac{3g}{2}))=-\frac{9}{2}g+mk+5m+6+m^{2}+\frac{9}{4}g^{2}-3mg.$
Let $f(x)=-\frac{9}{2}x+mk+5m+6+m^{2}+\frac{9}{4}x^{2}-3mx, 4\leq x\leq \frac{2m}{3}.$ Since $\frac{df(x)}{dx}=-\frac{9}{2}+\frac{9x}{2}-3m<0$, $f(x)$ is a strictly monotone decreasing function. So when $g$ is even, $M(C_{1,n,\frac{g}{2},\frac{g}{2},\frac{g}{2}}^{k}(m-\frac{3g}{2}))\leq M(C_{1,n,2,2,2}^{k}(m-6))$ for $4\leq g\leq \frac{2m}{3}$ with the equation if and only if $g=4$. Since the maximum degree of $ C_{3,n,\frac{g}{2}-1,\frac{g}{2}+1,\frac{g}{2}-1}^{k}$ is 2, $M(C_{1,n,2,2,2}^{k}(m-6))> M(C_{3,n,\frac{g}{2}-1,\frac{g}{2}+1,\frac{g}{2}-1}^{k})$ for $g=\frac{2m}{3}+\frac{2}{3}.$ %Note that when $3~|~m+1$, $C_{3,n,\frac{g}{2}-1,\frac{g}{2}+1,\frac{g}{2}-1}^{k}$ exists.
%Hence, $C_{1,n,2,2,2}^{k}(m-6)$ is the last hypergraph in an $S$-order of all linear bicyclic $k$-uniform hypergraphs in $C_{n,p,q,l}^{k}$.

When $g$ is odd, we have $M(C_{2,n,\lfloor\frac{g}{2}\rfloor,\lceil\frac{g}{2}\rceil,\lfloor\frac{g}{2}\rfloor}^{k}(m-g-\lfloor\frac{g}{2}\rfloor))=-6g+\frac{23}{4}+mk+6m+m^{2}+\frac{9g^{2}}{4}-3mg$.
%\begin{align*}
%&M(C_{2,n,\lfloor\frac{g}{2}\rfloor,\lceil\frac{g}{2}\rceil,\lfloor\frac{g}{2}\rfloor}^{k}(m-g-\lfloor\frac{g}{2}\rfloor))\\
%&=(g+\frac{g-1}{2}-1)(k-2)+(k-3)+(m-g-\frac{g-1}{2})(k-1)+4(g+\frac{g-1}{2}-1)\\
%&+(3+m-g-\frac{g-1}{2})^{2}\\
%&=-6g+\frac{23}{4}+mk+6m+m^{2}+\frac{9g^{2}}{4}-3mg.
%\end{align*}
%Let $f(x)=-6x+\frac{23}{4}+mk+6m+m^{2}+\frac{9x^{2}}{4}-3mx, 3\leq x\leq \frac{2m+1}{3}.$ Since $\frac{df(x)}{dx}=-6+\frac{9x}{2}-3m\leq-\frac{9}{2}<0$, $f(x)$ is a strictly monotone decreasing function. So $M(C_{2,n,\lfloor\frac{g}{2}\rfloor,\lceil\frac{g}{2}\rceil,\lfloor\frac{g}{2}\rfloor}^{k}(m-g-\lfloor\frac{g}{2}\rfloor))\leq M(C_{2,n,1,2,1}^{k}(m-4))$ for $3\leq g\leq \frac{2m+1}{3}$ with the equation if and only if $g=3$.
Let $f(x)=-6x+\frac{23}{4}+mk+6m+m^{2}+\frac{9x^{2}}{4}-3mx, 3\leq x\leq \frac{2m}{3}+\frac{1}{3}.$ Since $\frac{df(x)}{dx}=-6+\frac{9x}{2}-3m<0$, $f(x)$ is a strictly monotone decreasing function. So
when $g$ is odd, $M(C_{2,n,\lfloor\frac{g}{2}\rfloor,\lceil\frac{g}{2}\rceil,\lfloor\frac{g}{2}\rfloor}^{k}(m-g-\lfloor\frac{g}{2}\rfloor))\leq M(C_{2,n,1,2,1}^{k}(m-4))$ for $3\leq g\leq \frac{2m}{3}$ with the equation if and only if $g=3$. And $M(C_{2,n,1,2,1}^{k}(m-4))>M(C_{2,n,\lfloor\frac{g}{2}\rfloor,\lceil\frac{g}{2}\rceil,\lfloor\frac{g}{2}\rfloor}^{k})$ for $g=\frac{2m}{3}+\frac{1}{3}.$
Since the maximum degree of $C_{3,n,\frac{1}{2}g-\frac{3}{2},\frac{1}{2}g+\frac{3}{2},\frac{1}{2}g-\frac{1}{2}}^{k}$ is 2, $M(C_{1,n,2,2,2}^{k}(m-6))> M(C_{3,n,\frac{1}{2}g-\frac{3}{2},\frac{1}{2}g+\frac{3}{2},\frac{1}{2}g-\frac{1}{2}}^{k})$ for $g=\frac{2m}{3}+\frac{1}{3}.$


When $m\geq6$,
since $M(C_{1,n,2,2,2}^{k}(m-6))-M(C_{2,n,1,2,1}^{k}(m-4))=16-4m<0$, $C_{2,n,1,2,1}^{k}(m-4)$ is the hypergraph with maximum Zagreb index among all hypergraphs in $\mathcal{C}_{n}^{k}$.
%When $m=5$, obviously, $C_{2,n,1,2,1}^{k}(m-4)$ is the hypergraph with maximum Zagreb index among all hypergraphs in $\mathcal{C}_{n}^{k}$. Hence, for $m\geq 5$, $C_{2,n,1,2,1}^{k}(m-4)$ is the hypergraph with maximum Zagreb index among all hypergraphs in $\mathcal{C}_{n}^{k}$.
\end{proof}



The following Theorem gives the hypergraph with maximum Zagreb index among all linear bicyclic $k$-uniform hypergraphs with $n$ vertices, $m$ edges and girth $g$. And gives the hypergraph with maximum Zagreb index among all linear bicyclic $k$-uniform hypergraphs with $n$ vertices and $m$ edges.
\begin{thm}
%For $m\geq3$, in an $S$-order of $\textbf{U}_{g,f}$ the last hypergraph is $F_{g,f}$.%����ֻ��һ��Ȧ
For $k\geq3$ and $m\geq 2g$, when $g$ is even, $C_{1,n,\frac{g}{2},\frac{g}{2},\frac{g}{2}}^{k}(m-\frac{3g}{2})$ is the hypergraph with maximum Zagreb index among all linear bicyclic $k$-uniform hypergraphs with $n$ vertices, $m$ edges and girth $g$.
When $g$ is odd, $C_{2,n,\lfloor\frac{g}{2}\rfloor,\lceil\frac{g}{2}\rceil,\lfloor\frac{g}{2}\rfloor}^{k}(m-g-\lfloor\frac{g}{2}\rfloor)$ is the hypergraph with maximum Zagreb index among all linear bicyclic $k$-uniform hypergraphs with $n$ vertices, $m$ edges and girth $g$.

For $m\geq 6$, $C_{2,n,1,2,1}^{k}(m-4)$ is the hypergraph with maximum Zagreb index among all linear bicyclic $k$-uniform hypergraphs with $n$ vertices and $m$ edges.
\end{thm}
\begin{proof}
When $g$ is even,
$M(C_{1,n,\frac{g}{2},\frac{g}{2},\frac{g}{2}}^{k}(m-\frac{3g}{2}))-M(B_{1,n,g,0,g}^{k}(m-2g))=\frac{11}{2}g-2m-2-\frac{7}{4}g^{2}+mg$, $4\leq g\leq \frac{m}{2}$.
Let $f(x)=\frac{11}{2}x-2m-2-\frac{7}{4}x^{2}+mx$. The roots of $f(x)$ are easily obtained as $x_{1}=\frac{11+2m-2\sqrt{(m-\frac{3}{2})^{2}+14}}{7}$ and $x_{2}=\frac{11+2m+2\sqrt{(m-\frac{3}{2})^{2}+14}}{7}$. Since
$x_{1}<\frac{11+2m-2(m-\frac{3}{2})}{7}=2<4$ and $x_{2}>\frac{11+2m+2(m-\frac{3}{2})}{7}=\frac{4m+8}{7}>\frac{m}{2}$, when $4\leq g\leq \frac{m}{2}$, $M(C_{1,n,\frac{g}{2},\frac{g}{2},\frac{g}{2}}^{k}(m-\frac{3g}{2}))-M(B_{1,n,g,0,g}^{k}(m-2g))>0.$
%Since $\frac{d^{2}f(x)}{dx^{2}}<0$, $f(4)>0$ and  $f(\frac{m}{2})>0$, when $4\leq g\leq \frac{m}{2}$, $M(C_{1,n,\frac{g}{2},\frac{g}{2},\frac{g}{2}}^{k}(m-\frac{3g}{2}))-M(B_{1,n,g,0,g}^{k}(m-2g))>0.$
Therefore, when $g$ is even, $C_{1,n,\frac{g}{2},\frac{g}{2},\frac{g}{2}}^{k}(m-\frac{3g}{2})$ is the hypergraph with maximum Zagreb index among all linear bicyclic $k$-uniform hypergraphs with $n$ vertices, $m$ edges and girth $g$.

When $g$ is odd,
$M(C_{2,n,\lfloor\frac{g}{2}\rfloor,\lceil\frac{g}{2}\rceil,\lfloor\frac{g}{2}\rfloor}^{k}(m-g-\lfloor\frac{g}{2}\rfloor))-M(B_{1,n,g,0,g}^{k}(m-2g))=4g-m-\frac{9}{4}-\frac{7}{4}g^{2}+mg, 3\leq g\leq \frac{m}{2}$.
Let $f(x)=4x-m-\frac{9}{4}-\frac{7}{4}x^{2}+mx$. 
The roots of $f(x)$ are easily obtained as $x_{1}=\frac{8+2m-2\sqrt{(m+\frac{1}{2})^{2}}}{7}=1$ and $x_{2}=\frac{8+2m+2\sqrt{(m+\frac{1}{2})^{2}}}{7}=\frac{4m+9}{7}$. Obviously,
$x_{1}<3, x_{2}>\frac{m}{2}.$ So, when $3\leq g\leq \frac{m}{2}$, $M(C_{2,n,\lfloor\frac{g}{2}\rfloor,\lceil\frac{g}{2}\rceil,\lfloor\frac{g}{2}\rfloor}^{k}(m-g-\lfloor\frac{g}{2}\rfloor))-M(B_{1,n,g,0,g}^{k}(m-2g))>0.$
%Since $\frac{d^{2}f(x)}{dx^{2}}<0$, $f(3)>0$ and  $f(\frac{m}{2})>0$, when $3\leq g\leq \frac{m}{2}$, $M(C_{2,n,\lfloor\frac{g}{2}\rfloor,\lceil\frac{g}{2}\rceil,\lfloor\frac{g}{2}\rfloor}^{k}(m-g-\lfloor\frac{g}{2}\rfloor))-M(B_{1,n,g,0,g}^{k}(m-2g))>0.$
%When $g$ is odd,
%$M(C_{1,n,\lfloor\frac{g}{2}\rfloor,\lceil\frac{g}{2}\rceil,\lceil\frac{g}{2}\rceil}^{k}(m-g-\lceil\frac{g}{2}\rceil)-M(B_{1,n,g,0,g}^{k}(m-2g))=7g-3m-\frac{13}{4}-\frac{7}{4}g^{2}+mg.$
%Let $f(x)=7x-3m-\frac{13}{4}-\frac{7}{4}x^{2}+mx$. Then $x_{1}=\frac{14+2m-2\sqrt{(m-\frac{7}{2})^{2}+14}}{7}$ or $x_{2}=\frac{14+2m+2\sqrt{(m-\frac{7}{2})^{2}+14}}{7}$.
%$x_{1}<\frac{14+2m-2(m-\frac{7}{2})}{7}=3, x_{2}>\frac{14+2m+2(m-\frac{7}{2})}{7}=\frac{4m+7}{7}>\frac{m}{2}.$ So, when $3\leq g\leq \frac{m}{2}$, $M(C_{1,n,\lfloor\frac{g}{2}\rfloor,\lceil\frac{g}{2}\rceil,\lceil\frac{g}{2}\rceil}^{k}(m-g-\lceil\frac{g}{2}\rceil)-M(B_{1,n,g,0,g}^{k}(m-2g))>0.$
Hence, when $g$ is odd, $C_{2,n,\lfloor\frac{g}{2}\rfloor,\lceil\frac{g}{2}\rceil,\lfloor\frac{g}{2}\rfloor}^{k}(m-g-\lfloor\frac{g}{2}\rfloor)$ is the hypergraph with maximum Zagreb index among all linear bicyclic $k$-uniform hypergraphs with $n$ vertices, $m$ edges and girth $g$.

%Therefore, the hypergraph with maximum Zagreb index among all linear bicyclic $k$-uniform hypergraphs with $n$ vertices and $m$ edges is in $\{C_{1,n,\frac{g}{2},\frac{g}{2},\frac{g}{2}}^{k}(m-\frac{3g}{2})~|~g~\text{is even,~}4\leq g\leq \frac{2m}{3}}\bigcup \{C_{2,n,\lfloor\frac{g}{2}\rfloor,\lceil\frac{g}{2}\rceil,\lfloor\frac{g}{2}\rfloor}^{k}(m-g-\lfloor\frac{g}{2}\rfloor)~|~ g~ \text{is odd,~}3\leq g\leq \frac{2m}{3}\}$.
From the proof of Theorem \ref{q1}, %we know $M(C_{1,n,\frac{g}{2},\frac{g}{2},\frac{g}{2}}^{k}(m-\frac{3g}{2}))\leq M(C_{1,n,2,2,2}^{k}(m-6))$ for $4\leq g\leq \frac{2m}{3}$ with the equation if and only if $g=4$. $M(C_{2,n,\lfloor\frac{g}{2}\rfloor,\lceil\frac{g}{2}\rceil,\lfloor\frac{g}{2}\rfloor}^{k}(m-g-\lfloor\frac{g}{2}\rfloor))\leq M(C_{2,n,1,2,1}^{k}(m-4))$ for $3\leq g\leq \frac{2m+1}{3}$ with the equation if and only if $g=3$.
$C_{2,n,1,2,1}^{k}(m-4)$ is the hypergraph with maximum Zagreb index among all linear bicyclic $k$-uniform hypergraphs with $n$ vertices and $m$ edges.
\end{proof}


\vspace{3mm}

\noindent
\textbf{Acknowledgments}
\vspace{3mm}
\noindent

%The authors would like to thank the reviewers for giving valuable suggestions.
This work is supported by the National Natural Science Foundation of China (No. 12071097, No. 12371344), the Natural Science Foundation for The Excellent Youth Scholars of the Heilongjiang Province (No. YQ2022A002) and the Fundamental Research Funds for the
Central Universities.
%Heilongjiang Province (No. QC2018002) and the Fundamental Research Funds for
%the Central Universities.







\section*{References}
\bibliographystyle{unsrt}
\bibliography{pbib}
\end{spacing}
\end{document}
%%
%% End of file `elsarticle-template-1-num.tex'.
